% Preamble template for Tufte-style applied econometrics lecture notes.
% Do not compile this file directly; compile the main notes file instead.
\documentclass{tufte-handout}

\usepackage[T1]{fontenc}
\usepackage[utf8]{inputenc}
\usepackage{ntheorem}
\usepackage{graphicx}
\usepackage{amsmath}
\usepackage{amssymb}
\usepackage{mathtools}
\usepackage{hyperref}
\usepackage{epigraph}
\usepackage{booktabs}
\usepackage{array}
\theoremstyle{break}

\hypersetup{
  colorlinks,
  urlcolor = blue,
  pdfauthor={Swapnil Singh},
  pdfkeywords={econometrics, causal inference, applied econometrics},
  pdftitle={Lecture Notes for Applied Econometrics},
  pdfpagemode=UseNone
}

\newtheorem{ruleN}{Rule}
\newtheorem{thmN}{Theorem}
\newtheorem{assN}{Assumption}
\newtheorem{defN}{Definition}
\newtheorem{exmp}{Example}
\newtheorem{cmt}{Comment}
\newtheorem{proof}{Proof}

\newcommand{\continuation}{??}
\newtheorem*{excont}{Example \continuation}
\newenvironment{continueexample}[1]
 {\renewcommand{\continuation}{\ref{#1}}\excont[continued]}
 {\endexcont}

\newcommand{\bY}{\mathbf{Y}}
\newcommand{\bX}{\mathbf{X}}
\newcommand{\bD}{\mathbf{D}}
\newcommand{\E}{\mathbb{E}}
\newcommand{\Prob}{\mathbb{P}}
\newcommand{\Var}{\mathrm{Var}}
\newcommand{\Cov}{\mathrm{Cov}}
\newcommand{\se}{\mathrm{SE}}
\newcommand{\CI}{\mathrm{CI}}
\newcommand{\ATE}{\mathrm{ATE}}
\newcommand{\ATT}{\mathrm{ATT}}
\newcommand{\TOT}{\mathrm{TOT}}
\newcommand{\LATE}{\mathrm{LATE}}
\newcommand{\ITT}{\mathrm{ITT}}
\newcommand{\RD}{\mathrm{RD}}
\newcommand{\DiD}{\mathrm{DiD}}
\newcommand{\MMSE}{\mathrm{MMSE}}
\newcommand{\OLS}{\mathrm{OLS}}
\newcommand{\Normal}{\mathcal{N}}
\newcommand{\plim}{\operatorname*{plim}}
\DeclareMathOperator*{\argmin}{arg\,min}
\DeclareMathOperator{\Supp}{Supp}

\newcommand\independent{\protect\mathpalette{\protect\independenT}{\perp}}
\def\independenT#1#2{\mathrel{\rlap{$#1#2$}\mkern2mu{#1#2}}}
\newcommand{\indep}{\perp\!\!\!\perp}

\usepackage{cleveref}
\crefname{appsec}{appendix}{appendices}
\crefname{appsubsec}{appendix}{appendices}
\crefname{assumption}{assumption}{assumptions}
\crefname{equation}{equation}{equations}
\crefname{exmp}{example}{examples}
\crefname{assN}{assumption}{assumptions}
\crefname{cmt}{comment}{comments}
\crefname{defN}{definition}{definitions}

\usepackage[nolist]{acronym}
\begin{acronym}
  \acro{ATE}{average treatment effect}
  \acro{ATT}{average treatment effect on the treated}
  \acro{CIA}{conditional independence assumption}
  \acro{CI}{confidence interval}
  \acro{CLT}{central limit theorem}
  \acro{DiD}{difference-in-differences}
  \acro{FE}{fixed effects}
  \acro{FWL}{Frisch-Waugh-Lovell}
  \acro{IV}{instrumental variables}
  \acro{LATE}{local average treatment effect}
  \acro{OLS}{ordinary least squares}
  \acro{OVB}{omitted-variable bias}
  \acro{PO}{potential outcomes}
  \acro{RCT}{randomized controlled trial}
  \acro{RD}{regression discontinuity}
  \acro{SUTVA}{stable unit treatment value assignment}
  \acro{TWFE}{two-way fixed effects}
\end{acronym}

\def\inprobHIGH{\,{\buildrel p \over \rightarrow}\,}
\def\inprob{\,{\inprobHIGH}\,}
\def\indistHIGH{\,{\buildrel d \over \rightarrow}\,}
\def\indist{\,{\indistHIGH}\,}

\usepackage[many]{tcolorbox}
\definecolor{main}{HTML}{3F6EA7}
\definecolor{sub}{HTML}{F2F7FD}
\definecolor{sub2}{HTML}{FFF5E8}
\definecolor{mainline}{HTML}{D3DFEE}
\definecolor{warn}{HTML}{C7771B}
\definecolor{warnline}{HTML}{EAD9C0}

\tcbset{
    enhanced,
    breakable,
    colback = white,
    boxrule = 0.5pt,
    arc = 2pt,
    left = 9pt,
    right = 9pt,
    top = 7pt,
    bottom = 7pt,
    before skip = 0.3cm,
    after skip = 0.45cm
}

\newtcolorbox{boxD}{
    colback = sub,
    colframe = mainline,
    borderline west = {2.2pt}{0pt}{main},
    borderline north = {0.6pt}{0pt}{main!25},
    borderline south = {0.6pt}{0pt}{main!15}
}

\newtcolorbox{boxF}{
    colback = sub2,
    colframe = warnline,
    borderline west = {2.2pt}{0pt}{warn},
    borderline north = {0.6pt}{0pt}{warn!25},
    borderline south = {0.6pt}{0pt}{warn!15}
}

\usepackage{colortbl}
\usepackage{tikz}
\usepackage{verbatim}
\usetikzlibrary{positioning}
\usetikzlibrary{snakes}
\usetikzlibrary{calc}
\usetikzlibrary{arrows}
\usetikzlibrary{decorations.markings}
\usetikzlibrary{shapes.misc}
\usetikzlibrary{matrix,shapes,arrows,fit,tikzmark}


\title{Applied Econometrics: Session 4 Notes}
\author{Regression, Dummy Variables, Grouped Data, and Controls}
\date{}

\begin{document}

\maketitle

\begin{abstract}
This handout explains why linear regression is so useful in applied causal
work, and why it is also easy to misuse. We connect regression to potential
outcomes, conditional expectation functions, dummy variables, grouped data,
and the choice of controls. The main lesson is practical: a regression
coefficient is credible only when the comparison it represents is credible.
\end{abstract}

\begin{boxD}
\textbf{Reading and objective.} These notes accompany Applied Econometrics
Session 4 and are based on chapters 05, 06, and 07 of the Python Causality
Handbook: \emph{The Unreasonable Effectiveness of Linear Regression},
\emph{Grouped and Dummy Regression}, and \emph{Beyond Confounders}.

By the end of the session you should be able to explain what regression is
comparing, derive the omitted-variable-bias formula, use dummy variables and
grouped regressions correctly, and decide whether a proposed control is a
confounder, a precision control, a treatment predictor, a mediator, or a
collider.
\end{boxD}

\section{The Big Picture}

Regression is often introduced as a prediction tool. In this course we use it
mostly as a comparison tool. A regression with a treatment dummy asks whether
treated units have different outcomes from untreated units. A regression with
controls asks whether treated and untreated units have different outcomes
\emph{after comparing units with the same or similar controls}. Dummy
variables let those comparisons be flexible across groups. Grouped regression
shows that many regressions are just weighted averages of group-level
comparisons.

The danger is that the software always returns a coefficient. It does not know
whether the coefficient is causal. That judgment comes from the research
design.

\begin{boxF}
\textbf{The habit to build.} Every time you see a regression coefficient, ask:
``Which units are being compared? Why are those units comparable? What has
been held fixed? What has been accidentally held fixed?''
\end{boxF}

\section{Potential Outcomes Reminder}

Let $D_i \in \{0,1\}$ denote treatment. Unit $i$ has two potential outcomes:
\[
  Y_i(1) \quad \text{if treated}, \qquad
  Y_i(0) \quad \text{if untreated}.
\]
We observe only
\[
  Y_i = D_iY_i(1) + (1-D_i)Y_i(0).
\]
The individual treatment effect is $Y_i(1)-Y_i(0)$, but for any one unit one
of the two potential outcomes is missing. The average treatment effect is
\[
  ATE = \E[Y_i(1)-Y_i(0)].
\]
For the treated units, the target is often
\[
  ATT = \E[Y_i(1)-Y_i(0)\mid D_i=1].
\]

\begin{defN}[Selection bias]
The raw treated-control difference can be decomposed as
\[
\begin{aligned}
  \E[Y_i\mid D_i=1]-\E[Y_i\mid D_i=0]
  &= \E[Y_i(1)\mid D_i=1]-\E[Y_i(0)\mid D_i=0] \\
  &= ATT
  + \left\{\E[Y_i(0)\mid D_i=1]-\E[Y_i(0)\mid D_i=0]\right\}.
\end{aligned}
\]
The second term is selection bias: the difference that would have existed even
without treatment.
\end{defN}

This formula is the reason we need a research design. Treated people often
differ from untreated people before treatment. Students with more education
may have higher wages not only because education raises wages, but also
because they differ in ability, family background, local labor markets, or
motivation.

\section{Treatment Dummy Regression}

Start with the simplest regression:
\[
  Y_i = \alpha + \tau D_i + u_i.
\]
The fitted value for an untreated unit is $\hat\alpha$. The fitted value for a
treated unit is $\hat\alpha+\hat\tau$. Therefore
\[
  \hat\tau
  =
  \bar Y_1-\bar Y_0,
\]
where $\bar Y_1$ is the sample mean among treated units and $\bar Y_0$ is the
sample mean among untreated units.

\begin{thmN}[Dummy regression equals a difference in means]
In an OLS regression of $Y_i$ on a constant and a binary variable $D_i$, the
coefficient on $D_i$ equals the treated sample mean minus the untreated sample
mean.
\end{thmN}

\begin{proof}
OLS minimizes
\[
  \sum_{i:D_i=0}(Y_i-a)^2+\sum_{i:D_i=1}(Y_i-a-b)^2.
\]
For untreated units the fitted value is $a$; for treated units it is $a+b$.
The least-squares fitted value within each group is the group mean, so
$\hat a=\bar Y_0$ and $\hat a+\hat b=\bar Y_1$. Thus
$\hat b=\bar Y_1-\bar Y_0$.
\end{proof}

\begin{exmp}[Job training]
Suppose the average earnings of trained workers are 24,000 euros and the
average earnings of untrained workers are 20,000 euros. The treatment-dummy
regression coefficient is 4,000. This is a causal estimate only if the trained
and untrained workers would have had the same average earnings without the
program. If better workers selected into training, the coefficient overstates
the effect.
\end{exmp}

\section{Conditional Expectation Functions}

The conditional expectation function, or CEF, is
\[
  m(x)=\E[Y_i\mid X_i=x].
\]
It is the best mean-squared-error predictor of $Y$ using $X$. In causal work,
the CEF is useful because it describes how average outcomes vary with
observables. It is not automatically causal. If $X$ includes education, then
$\E[wage\mid education=16]$ is the average wage among people with 16 years of
education. It is not automatically the wage a randomly chosen person would
earn if forced to obtain 16 years of education.

\begin{defN}[Conditional independence]
A control vector $X_i$ solves selection on observables if
\[
  \{Y_i(1),Y_i(0)\}\perp D_i \mid X_i.
\]
This says that after conditioning on $X_i$, treatment assignment is as good as
random.
\end{defN}

Under conditional independence and overlap, treatment-control comparisons
within values of $X$ are causal:
\[
  \E[Y_i\mid D_i=1,X_i=x]-\E[Y_i\mid D_i=0,X_i=x]
  =
  \E[Y_i(1)-Y_i(0)\mid X_i=x].
\]
Regression then becomes a way of averaging many conditional comparisons.

\section{Population and Sample OLS}

Population OLS chooses coefficients that minimize expected squared error:
\[
  \beta^{pop}
  =
  \arg\min_b \E[(Y_i-X_i'b)^2].
\]
The first-order condition is
\[
  \E[X_i(Y_i-X_i'\beta^{pop})]=0.
\]
If $\E[X_iX_i']$ is invertible, then
\[
  \beta^{pop}=\E[X_iX_i']^{-1}\E[X_iY_i].
\]

Sample OLS replaces expectations by averages. With matrix $X$ and outcome
vector $Y$,
\[
  \hat\beta = (X'X)^{-1}X'Y,
\]
provided $X'X$ is invertible. The residual vector $\hat u=Y-X\hat\beta$
satisfies
\[
  X'\hat u=0.
\]
So each included regressor is uncorrelated with the OLS residual in the sample.
This is a mechanical property of OLS. It is not an assumption and it is not
evidence that the model is causal.

\section{Regression Anatomy and FWL}

Suppose we estimate
\[
  Y_i=\alpha+\tau D_i+X_i'\gamma+u_i.
\]
The Frisch-Waugh-Lovell theorem tells us exactly what $\hat\tau$ is doing.

\begin{thmN}[Frisch-Waugh-Lovell]
The coefficient $\hat\tau$ from the regression of $Y$ on a constant, $D$, and
$X$ is equal to the coefficient from a regression of residualized $Y$ on
residualized $D$, where both residualizations are from regressions on the
constant and $X$.
\end{thmN}

Let
\[
  \tilde D_i = D_i-\widehat{\E}[D_i\mid X_i],
  \qquad
  \tilde Y_i = Y_i-\widehat{\E}[Y_i\mid X_i].
\]
Then
\[
  \hat\tau
  =
  \frac{\sum_i \tilde D_i\tilde Y_i}{\sum_i \tilde D_i^2}.
\]
Thus multiple regression first removes from treatment the part predictable
from controls. Then it asks whether the remaining treatment variation predicts
the remaining outcome variation.

\begin{boxD}
\textbf{Plain-language FWL.} The coefficient on treatment comes from comparing
treated and untreated observations that differ in treatment after accounting
for the controls. If the controls almost perfectly predict treatment, there is
very little residual treatment variation left, and the estimate becomes noisy.
\end{boxD}

\section{When Regression Adjustment Is Causal}

Regression adjustment is credible when the conditional comparisons are
credible. A useful textbook condition is:
\[
  \{Y_i(1),Y_i(0)\}\perp D_i\mid X_i
  \quad \text{and} \quad
  0<P(D_i=1\mid X_i)<1.
\]
The first part is conditional independence. The second part is overlap:
for the same values of $X$, we need both treated and untreated units.

\begin{assN}[No bad controls]
The controls $X_i$ should be measured before treatment or otherwise not be
caused by treatment, and they should not be colliders created by conditioning
on common effects.
\end{assN}

\begin{assN}[Functional-form adequacy]
If treatment effects or baseline outcomes vary strongly with $X_i$, the
regression specification must be flexible enough to approximate the relevant
conditional means. Linearity is convenient, not magic.
\end{assN}

For weak students, the best way to think about this is blocked experiments.
If we compare treated and untreated men with similar age, education, and region,
we are pretending that treatment was randomized inside those blocks. Regression
does this blocking smoothly, with algebra rather than by manually forming
cells.

\section{Omitted-Variable Bias}

Suppose the true causal equation is
\[
  Y_i=\alpha+\tau D_i+\gamma Z_i+\varepsilon_i,
  \qquad
  \E[\varepsilon_i\mid D_i,Z_i]=0,
\]
but we estimate the short regression
\[
  Y_i=a+bD_i+e_i.
\]
Then the population short-regression coefficient is
\[
  b=\tau+\gamma\delta,
\]
where $\delta$ is the coefficient from the population regression of $Z_i$ on
$D_i$:
\[
  Z_i=\pi+\delta D_i+v_i.
\]

\begin{thmN}[Omitted-variable bias formula]
With one omitted variable $Z_i$ and one included treatment $D_i$,
\[
  \operatorname{Bias}(b)=b-\tau
  =
  \gamma
  \frac{\operatorname{Cov}(D_i,Z_i)}{\operatorname{Var}(D_i)}.
\]
\end{thmN}

\begin{proof}
The short-regression coefficient is
\[
  b=\frac{\operatorname{Cov}(D_i,Y_i)}{\operatorname{Var}(D_i)}.
\]
Substitute the true model:
\[
\begin{aligned}
  b
  &=\frac{\operatorname{Cov}(D_i,\alpha+\tau D_i+\gamma Z_i+\varepsilon_i)}
          {\operatorname{Var}(D_i)} \\
  &=\tau+\gamma
      \frac{\operatorname{Cov}(D_i,Z_i)}{\operatorname{Var}(D_i)}
      +\frac{\operatorname{Cov}(D_i,\varepsilon_i)}{\operatorname{Var}(D_i)}.
\end{aligned}
\]
The last term is zero by the true exogeneity condition, giving the result.
\end{proof}

\begin{boxF}
\textbf{Sign intuition.} Bias is positive if the omitted variable raises the
outcome and treated units have more of it. Bias is negative if the omitted
variable raises the outcome and treated units have less of it. Always write
the two signs: effect of omitted variable on $Y$, and association of omitted
variable with $D$.
\end{boxF}

\begin{exmp}[Education and wages]
Let $D_i$ be years of education and $Z_i$ be ability. If ability raises wages
($\gamma>0$) and high-ability people obtain more education
($\operatorname{Cov}(D_i,Z_i)>0$), then a wage regression omitting ability
overstates the return to education. If a policy forces low-ability students to
stay in school longer, the sign of the association could differ in that sample.
The bias formula is about the population being studied.
\end{exmp}

\section{Dummy Variables}

A dummy variable is a variable that takes values zero or one. Dummies are used
for binary treatments, categories, fixed effects, and flexible group means.

For a binary dummy $G_i$,
\[
  Y_i=\alpha+\beta G_i+u_i
\]
has fitted values
\[
  \hat Y_i =
  \begin{cases}
    \hat\alpha, & G_i=0,\\
    \hat\alpha+\hat\beta, & G_i=1.
  \end{cases}
\]
So $\hat\beta$ is the difference between group means.

\subsection{The Dummy-Variable Trap}

Suppose education has three categories: high school, college, graduate. If the
regression includes a constant and all three category dummies,
\[
  1 = 1\{HS\}+1\{College\}+1\{Graduate\}
\]
for every observation. One column is an exact linear combination of the
others. Then $X'X$ is not invertible.

The usual solution is to omit one category. That omitted category becomes the
reference group:
\[
  Y_i=\alpha+\beta_1 1\{College_i\}+\beta_2 1\{Graduate_i\}+u_i.
\]
Here $\alpha$ is the high-school mean, $\beta_1$ is college minus high school,
and $\beta_2$ is graduate minus high school.

\begin{boxF}
\textbf{Common mistake.} Students often say the omitted group is ``ignored.''
It is not ignored. It is the benchmark. Every included dummy coefficient is
interpreted relative to it.
\end{boxF}

\subsection{Saturated Dummies}

A saturated dummy regression includes one dummy for every group, except the
reference category if there is a constant. With only group dummies, it
reproduces the group means exactly. If group $g$ has mean $\bar Y_g$, the
fitted value for all observations in group $g$ is $\bar Y_g$.

This is why dummy regression is flexible. It does not force a straight-line
relationship between a categorical variable and the outcome. If wage rises
from 12 to 16 years of education and then flattens, education dummies can show
that pattern. A single linear education coefficient cannot.

\section{Grouped Regression and Weights}

Sometimes data are available only as cell means. Suppose each group $g$ has
mean outcome $\bar Y_g$, mean regressors $\bar X_g$, and cell size $n_g$.
Individual-level OLS solves
\[
  \min_b \sum_i (Y_i-X_i'b)^2.
\]
If all regressors are constant within group, the same coefficient can be
obtained from weighted grouped OLS:
\[
  \min_b \sum_g n_g(\bar Y_g-\bar X_g'b)^2.
\]

\begin{thmN}[Why grouped weights matter]
When grouped observations represent different numbers of individual
observations, the group-level regression should weight each group by its cell
size if the target is the individual-level regression.
\end{thmN}

Unweighted grouped regression gives each group the same importance. That may
be the right estimand if each group is itself the unit of interest. It is not
the same as an individual-level regression when one group contains 10 people
and another contains 10,000.

\begin{exmp}[Education cells]
Imagine wage data grouped by years of education. There may be many people with
12 years of education and few people with 20 years. If we regress cell mean
wages on education without weights, the small 20-year cell receives the same
weight as the large 12-year cell. Weighted grouped OLS gives each cell weight
proportional to the number of people in it.
\end{exmp}

\section{Interactions}

Interactions allow the treatment effect to differ across groups. With a binary
treatment $D_i$ and a binary group $G_i$,
\[
  Y_i=\alpha+\tau D_i+\eta G_i+\theta(D_iG_i)+u_i.
\]
The treatment effect for $G_i=0$ is $\tau$. The treatment effect for $G_i=1$
is $\tau+\theta$. Therefore $\theta$ is a difference in treatment effects
across groups.

\begin{boxD}
\textbf{Interpretation rule.} Do not interpret the coefficient on $D_i$ alone
as ``the treatment effect'' when the model contains $D_iG_i$. It is the
treatment effect for the reference group $G_i=0$.
\end{boxD}

Interactions are useful when theory suggests heterogeneity. A training program
may affect younger workers differently from older workers. A class-size
reduction may matter more in early grades. A tax change may affect liquidity
constrained households more than unconstrained households.

\section{Controls Beyond Confounders}

Not every control has the same role. The key control types are:
\[
\begin{array}{lll}
\text{Confounder} & \text{Causes } D \text{ and } Y & \text{Usually include} \\
\text{Outcome predictor} & \text{Predicts } Y \text{ only} & \text{Can improve precision} \\
\text{Treatment predictor} & \text{Predicts } D \text{ only} & \text{Can hurt precision} \\
\text{Mediator} & D \rightarrow M \rightarrow Y & \text{Do not include for total effect} \\
\text{Collider} & D \rightarrow C \leftarrow U \rightarrow Y & \text{Do not include}
\end{array}
\]

\subsection{Precision Controls}

In a randomized experiment, treatment is independent of potential outcomes.
Controls are not needed for identification. But controls that strongly predict
the outcome can reduce residual variance and improve precision:
\[
  Y_i=\alpha+\tau D_i+X_i'\gamma+u_i.
\]
If $X_i$ explains a large part of $Y_i$ but is unrelated to $D_i$ by
randomization, including $X_i$ can make $\hat\tau$ more precise without
changing its causal interpretation.

\subsection{Treatment Predictors}

A variable that predicts treatment but not the outcome can reduce useful
variation in treatment. In FWL language, it makes $\tilde D_i$ smaller without
making $\tilde Y_i$ much cleaner. Standard errors may rise. This is why
``include everything that predicts treatment'' is not a good rule.

\subsection{Mediators}

A mediator is caused by treatment and then affects the outcome:
\[
  D_i \rightarrow M_i \rightarrow Y_i.
\]
If the target is the total effect of treatment, do not control for $M_i$.
Controlling for the mediator blocks part of the treatment effect.

\begin{exmp}[Email reminder]
Suppose an email reminder increases payment partly because it makes people log
into an online account. If the outcome is payment and the mediator is login,
then controlling for login estimates the effect of the email holding login
fixed. That is not the total effect of the email.
\end{exmp}

\subsection{Colliders}

A collider is a common effect of two variables. Conditioning on it can create
a false association. A typical structure is
\[
  D_i \rightarrow C_i \leftarrow U_i \rightarrow Y_i.
\]
Even if $D_i$ was randomized, controlling for $C_i$ can open a path between
$D_i$ and unobserved $U_i$, biasing the coefficient on treatment.

\begin{boxF}
\textbf{Selection is often collider conditioning.} If we analyze only people
who applied for a program, only firms that survived, or only loans that were
approved, we may be conditioning on a variable affected by both treatment and
unobserved determinants of the outcome.
\end{boxF}

\section{Using the Python Example}

The companion script
\[
  \texttt{lectures/code/lecture-4-regression-controls.py}
\]
contains three classroom examples.

\begin{enumerate}
\item \textbf{Omitted-variable bias.} The script simulates education, wages,
latent ability, and family background. The short regression is biased because
education is correlated with omitted determinants of wages.
\item \textbf{Dummy and grouped regression.} The script compares a linear
education regression with a saturated education-dummy regression and shows
that weighted grouped OLS matches the individual-level regression when the
weights are cell sizes.
\item \textbf{Bad controls.} The script shows how adding a mediator changes
the estimand and how controlling for a collider can create bias even when
treatment was randomized.
\end{enumerate}

When using the script in class, ask students to predict the sign of each
coefficient before looking at the output. Prediction forces them to reason
about the comparison rather than memorize a formula.

\section{Practical Control-Selection Rules}

\begin{enumerate}
\item Start with the estimand. Are you trying to estimate a total effect, a
direct effect, an ATT, or a descriptive association?
\item Draw a causal story before opening the data. List variables that cause
treatment, cause the outcome, are caused by treatment, or are selection rules.
\item Include pre-treatment confounders when the identification strategy is
selection on observables.
\item Include strong pre-treatment outcome predictors when they improve
precision and do not damage overlap.
\item Be cautious with variables that mostly predict treatment. They may make
standard errors larger by removing treatment variation.
\item Do not control for mediators if the target is the total effect.
\item Do not control for colliders or post-treatment selection variables.
\item Check overlap. Regression cannot credibly compare treated units in
regions of $X$ where no controls exist.
\item Show robustness, but do not treat a long list of controls as proof of
causality.
\end{enumerate}

\section{Worked Problems}

\subsection{Problem 1: Treatment Dummy}

In a sample of 100 students, 40 attend tutoring and have average exam score
78. The 60 non-tutored students have average score 70. What is the coefficient
from regressing score on a constant and a tutoring dummy?

\textbf{Solution.} The coefficient is the treated mean minus the untreated
mean:
\[
  \hat\tau=78-70=8.
\]
The intercept is 70. The fitted score is 70 for non-tutored students and 78
for tutored students. This is causal only if tutored and non-tutored students
would have had the same average score without tutoring, or if assignment to
tutoring was random.

\subsection{Problem 2: OVB Sign}

You regress wages on education but omit parental income. Parental income
raises wages through networks and school quality. Students from richer
families obtain more education. What is the likely sign of omitted-variable
bias?

\textbf{Solution.} The omitted variable has a positive effect on wages:
$\gamma>0$. It is positively associated with education:
$\operatorname{Cov}(D,Z)>0$. Therefore the omitted-variable bias is positive.
The short regression likely overstates the causal return to education.

\subsection{Problem 3: Dummy Trap}

You include a constant, a male dummy, and a female dummy in a regression. Why
does the software drop one variable?

\textbf{Solution.} For every observation,
\[
  1=Male_i+Female_i.
\]
The constant is exactly the sum of the two dummies. The design matrix has
perfect multicollinearity, so $X'X$ is singular. Dropping one dummy creates a
reference group.

\subsection{Problem 4: Grouped Weights}

A grouped dataset has cell means by region. Region A has 1,000 people and
region B has 20 people. You want the individual-level relationship between
income and education. Should the grouped regression weight the regions
equally?

\textbf{Solution.} No. If the target is the individual-level regression,
region A should receive 1,000 units of weight and region B should receive 20.
Equal weights would estimate a region-level relationship where each region
counts the same.

\subsection{Problem 5: Mediator}

A training program increases certification, and certification raises wages.
You regress wages on training and certification. What does the training
coefficient estimate?

\textbf{Solution.} It estimates the effect of training holding certification
fixed. It does not estimate the total effect of training, because one channel
from training to wages has been blocked.

\section{Checklist for Students}

Before trusting a regression coefficient, check the following.

\begin{enumerate}
\item Can I state the estimand in words?
\item Can I write the regression equation and interpret every coefficient?
\item For a treatment coefficient, do I know which units are being compared?
\item Are the controls pre-treatment?
\item Are any controls mediators or colliders?
\item Is there overlap between treated and untreated units after controls?
\item If dummies are used, do I know the reference group?
\item If interactions are used, do I know which group the main effect refers
to?
\item If grouped data are used, are the weights aligned with the estimand?
\item Can I explain the likely direction of omitted-variable bias?
\end{enumerate}

\section{Practice Questions}

\begin{enumerate}
\item Prove that the coefficient on a binary treatment dummy equals the
difference in sample means.
\item Derive the omitted-variable-bias formula for one omitted variable.
\item Give an example of a variable that is a confounder in one study but a
mediator in another.
\item In the regression
\[
  Y_i=\alpha+\tau D_i+\eta Female_i+\theta(D_i\times Female_i)+u_i,
\]
what is the treatment effect for men? What is the treatment effect for women?
\item Explain why a saturated dummy regression produces group means.
\item A researcher has grouped data with cell sizes. When should she use
weighted grouped OLS? When might unweighted grouped OLS be the right object?
\item Why can an outcome predictor improve precision in a randomized
experiment?
\item Why can a treatment predictor increase the standard error of the
treatment coefficient?
\item Draw a simple DAG with a collider and explain why conditioning on the
collider creates bias.
\item Run \texttt{lectures/code/lecture-4-regression-controls.py}. For each
example, write one paragraph explaining what changed across specifications and
why.
\end{enumerate}

\end{document}
